\documentclass[11pt]{article}

\usepackage[utf8]{inputenc}
\usepackage[T1]{fontenc}
\usepackage{lmodern}
\usepackage{geometry}
\usepackage{amsmath,amssymb,amsfonts,amsthm,mathtools}
\usepackage{bm}
\usepackage{enumitem}
\usepackage{microtype}
\usepackage{hyperref}
\usepackage{titlesec}
\usepackage{booktabs}
\usepackage{array}
\usepackage{setspace}
\usepackage{mathrsfs}
\usepackage{physics}

\geometry{margin=1in}
\hypersetup{colorlinks=true, linkcolor=black, urlcolor=blue, citecolor=black}
\setlist[enumerate]{topsep=0.4em,itemsep=0.35em}
\setlist[itemize]{topsep=0.3em,itemsep=0.25em}
\titleformat{\section}{\large\bfseries}{\thesection.}{0.5em}{}
\titleformat{\subsection}{\normalsize\bfseries}{\thesubsection.}{0.5em}{}
\setstretch{1.06}

\newcommand{\Aether}{\AE{}ther}
\newcommand{\Aflow}{\AE{}ther-flow}
\newcommand{\TheoryName}{\Aflow{} Interpretation of Relativity}
\newcommand{\TheoryProgram}{\Aether{} / \Aflow{} framework}
\newcommand{\RespShapeReadout}{\widetilde q_{\mathrm{br}}}
\newcommand{\RespShapeCov}{\widetilde\kappa_{\mathrm{br}}}
\newcommand{\RespShapeRelax}{\widetilde\rho_{\mathrm{br}}}
\newcommand{\RespShapeWell}{\widetilde\lambda_{\mathrm{br}}}
\newcommand{\RespShapeEq}{\widetilde U_{\mathrm{br}}^{\ast}}
\newcommand{\RespRawReadout}{\widehat q_{\mathrm{br}}}
\newcommand{\RespRawRef}{\widehat U_{\mathrm{br}}^{\mathrm{ref}}}
\newcommand{\RespRawCov}{\widehat\kappa_{\mathrm{br}}}
\newcommand{\RespRawRelax}{\widehat\rho_{\mathrm{br}}}
\newcommand{\RespRawWell}{\widehat\lambda_{\mathrm{br}}}
\newcommand{\RespRawEq}{\widehat U_{\mathrm{br}}^{\ast}}
\newcommand{\RespVacPhase}{\phi_{\mathrm{br}}}
\newcommand{\CurvProfileCan}{F_{\mathrm{br}}^{\mathrm{can}}}
\newcommand{\DownPkg}{\mathscr P_{\mathrm{down}}}
\newcommand{\ShapeTuple}{\mathfrak S_{\mathrm{br}}}
\newcommand{\RawTuple}{\mathfrak R_{\mathrm{br}}}
\newcommand{\DeepSector}{\mathfrak O_{\mathrm{br}}}
\newcommand{\ScaleMod}{s_{\mathrm{br}}}

\newtheorem{definition}{Definition}
\newtheorem{proposition}{Proposition}
\newtheorem{remark}{Remark}
\newtheorem{corollary}{Corollary}
\newtheorem{theorem}{Theorem}

\title{\textbf{The \TheoryName{}}\\[0.4em]
\large Nonlocal Bridge Self-Response Off-Shell Profile-Selection Bridge-Order Orbit-Shape/Modulus-Origin Note}
\author{}
\date{}

\begin{document}

\maketitle

\begin{abstract}
The flagship exact-theory statement of \TheoryName{} has already crossed its threshold and is not at stake here. The active exact-closure sequence already supplies the public theory line. The present manuscript acts only on the optional deeper derivational bridge below that line.

The raw-vacuum-orbit-origin note narrowed the remaining burden sharply. The raw positive vacuum orbit is no longer terminal input: it is generated from a deeper primitive orbit-modulus sector while preserving the same density/current block, the same composite primitive bridge-order block, the same phase-neutral minimizing branch, the same canonical profile, and the same downstream dressed bridge map. That note therefore crossed the honest raw-orbit threshold already. The next move, if the derivational line continues, is not another same-layer repair.

This note records the exact still-deeper burden. First, it defines what a successful orbit-shape/modulus-origin continuation must accomplish below the raw-vacuum-orbit layer: it must derive or sharply constrain the unit-reference orbit-shape block and the positive orbit modulus from still more primitive substrate structure. Second, it proves that any such admissible deeper continuation must leave the same downstream chain unchanged. Third, it proves that the derived positive scale orbit and the residual global \(O(2)\) vacuum angle may be revisited only if that deeper layer supplies genuine quotient or locking structure. Otherwise the current verdicts remain in force: canonical normalization is still the honest representative choice for the positive scale orbit, and the global angle remains residual.

The present note is therefore a burden-fixing manuscript, not a claim that the deeper derivation has already been completed. Its positive result is architectural discipline: the raw-vacuum-orbit note is confirmed as a genuine threshold crossing, and the next honest derivational note is fixed precisely.
\end{abstract}

\section{Introduction}

The active repository no longer needs another derivational manuscript in order to justify its public theory line. The exact-closure sequence already presents \TheoryName{} as a disciplined exact relativistic theory statement built on the \Aether{} / \Aflow{} ontology. Any continuation of the deeper bridge is therefore optional foundational work rather than a prerequisite for the flagship exact theory claim.

That distinction matters for the present note. The raw-vacuum-orbit-origin manuscript did not leave the project stuck at the same conceptual layer. It moved the burden genuinely one step deeper by deriving the raw positive vacuum orbit from a primitive orbit-modulus sector while preserving the whole already-positive downstream package. The next honest continuation, if one is written, must therefore go below the raw orbit itself rather than attempt another same-layer repair.

\paragraph{Framework and claim boundary.}
\TheoryProgram{} denotes the broader \Aether{} / \Aflow{} framework built on the claims that the \Aether{} is the underlying four-dimensional substrate of reality, the \Aflow{} is its intrinsic ordered motion, observed three-dimensional space is the local experiential slice of that deeper substrate, S-time is the experienced order of change arising from matter, light, and the \Aflow{}, and observed expansion is the three-dimensional appearance of deeper four-dimensional ordered motion. Within that framework, \TheoryName{} denotes the exact-closure theory in which the effective gravitational dynamics are adopted to be exactly Einsteinian with universal matter coupling. In the active sequence, \emph{adoption} means use of that established relativistic dynamics without claiming substrate derivation, while \emph{derivation} is reserved for a first-principles recovery from explicit substrate variables.

This note stays strictly within that boundary. It does not claim that the still deeper orbit-shape/modulus origin has already been solved. It claims something narrower and still useful: the project now knows exactly what that next note must do, exactly what it must preserve, and exactly when the scale-orbit and vacuum-angle verdicts may honestly change.

\section{Current Endpoint of the Derivational Bridge}

\begin{definition}[Current deepest recorded orbit data]
\label{def:current}
At current scope, the deepest recorded bridge-order layer consists of:
\begin{enumerate}
    \item the unit-reference orbit-shape block
    \begin{equation}
    \ShapeTuple
    :=
    (\RespShapeReadout,\RespShapeCov,\RespShapeRelax,\RespShapeWell,\RespShapeEq),
    \label{eq:shapetuple}
    \end{equation}
    \item a positive orbit modulus
    \begin{equation}
    \ScaleMod>0,
    \label{eq:scalemod}
    \end{equation}
    \item the induced raw positive vacuum orbit
    \begin{equation}
    \RawTuple(\ScaleMod)
    :=
    (\RespRawReadout,\RespRawRef,\RespRawCov,\RespRawRelax,\RespRawWell,\RespRawEq)
    \label{eq:rawtuple}
    \end{equation}
    with
    \begin{equation}
    \RespRawReadout=\ScaleMod\,\RespShapeReadout,
    \qquad
    \RespRawRef=\ScaleMod,
    \qquad
    \RespRawCov=\ScaleMod^2\,\RespShapeCov,
    \label{eq:rawliftA}
    \end{equation}
    \begin{equation}
    \RespRawRelax=\ScaleMod^2\,\RespShapeRelax,
    \qquad
    \RespRawWell=\RespShapeWell,
    \qquad
    \RespRawEq=\ScaleMod\,\RespShapeEq,
    \label{eq:rawliftB}
    \end{equation}
    \item the already-recovered downstream package \(\DownPkg\), namely the same density/current block, the same composite primitive bridge-order block, the same phase-neutral minimizing branch, the same canonical profile
    \begin{equation}
    \CurvProfileCan(x)=1+\RespShapeEq\,x^2
    =
    1+\frac{\RespRawEq}{\RespRawRef}x^2,
    \label{eq:canonicalprofile}
    \end{equation}
    and the same dressed bridge map as in the current raw-vacuum-orbit continuation.
\end{enumerate}
\end{definition}

\begin{proposition}[The raw-vacuum-orbit layer has already crossed its threshold]
\label{prop:nosamelayer}
Relative to Definition \ref{def:current}, the next honest derivational burden is not another repair at the raw-orbit layer itself. It is deeper again: the origin, anchoring, or quotient of \(\ShapeTuple\) and \(\ScaleMod\).
\end{proposition}

\begin{proof}
The current raw-vacuum-orbit result already achieves the correct burden at its own layer: it generates the raw positive vacuum orbit from a deeper primitive orbit-modulus sector and proves that the already-positive downstream package \(\DownPkg\) survives unchanged. Therefore the raw-orbit slot is no longer the active incompleteness. Any next note that merely reparameterizes \(\RawTuple(\ScaleMod)\) without going below \(\ShapeTuple\) and \(\ScaleMod\) does not address the actual remaining gap. The remaining gap is precisely the status of the unit-reference shape data and the positive modulus that still feed the raw orbit.
\end{proof}

\begin{remark}
Proposition \ref{prop:nosamelayer} is architectural, not negative. It does not say the bridge program is blocked. It says the continuation direction is now disciplined: either consolidate the flagship exact-theory line, or go one layer deeper than the raw orbit itself.
\end{remark}

\section{What the Next Orbit-Shape/Modulus Note Must Accomplish}

\begin{definition}[Admissible deeper orbit-shape/modulus continuation]
\label{def:admissible}
A still deeper continuation below the current raw-vacuum-orbit layer is called \emph{admissible} if it introduces a more primitive substrate sector \(\DeepSector\) together with maps
\begin{equation}
\DeepSector
\longmapsto
(\ShapeTuple,\ScaleMod),
\label{eq:deepermap}
\end{equation}
such that all of the following hold.
\begin{enumerate}
    \item \textbf{Deeper origin or sharp constraint.} The note derives or sharply constrains the unit-reference shape block \(\ShapeTuple\) and the positive modulus \(\ScaleMod\) from \(\DeepSector\) rather than leaving them as unexplained terminal data.
    \item \textbf{Same raw-orbit lift.} The induced raw orbit is still given by \eqref{eq:rawliftA}--\eqref{eq:rawliftB}.
    \item \textbf{Same phase-neutral branch.} The minimizing vacuum branch remains phase-neutral, so the distinguished minimizing family still has constant \(\RespVacPhase\) and no new observer-accessible angular forcing.
    \item \textbf{Same downstream dependence.} Every downstream object in \(\DownPkg\) depends on the deeper layer only through the induced pair \((\ShapeTuple,\ScaleMod)\) and the same phase-neutral minimizing branch.
\end{enumerate}
\end{definition}

\begin{remark}
Definition \ref{def:admissible} leaves open whether the deeper sector determines \(\ShapeTuple\) uniquely, constrains it to a narrow admissible class, or derives it only modulo a later quotient. The present point is not to prejudge the mechanism. It is to fix the burden below the raw-orbit layer precisely enough that a future manuscript can be judged honestly.
\end{remark}

\begin{corollary}[The next note has three required tasks]
\label{cor:threetasks}
If the derivational line continues past the raw-vacuum-orbit note, the next honest manuscript must do exactly three things:
\begin{enumerate}
    \item derive or sharply constrain \(\ShapeTuple\) and \(\ScaleMod\) from still more primitive substrate structure;
    \item prove that the same downstream package \(\DownPkg\) survives unchanged;
    \item revisit the scale-orbit and vacuum-angle verdicts only if the deeper layer naturally supplies genuine quotient or locking structure.
\end{enumerate}
\end{corollary}

\begin{proof}
Item (1) is exactly the remaining burden identified in Proposition \ref{prop:nosamelayer}. Item (2) is required because the active line is already positive downstream, so a deeper continuation is admissible only if it preserves rather than rewrites those results. Item (3) follows because the current scale-orbit and vacuum-angle verdicts are already disciplined; they should change only in response to genuinely new deeper structure, not by assertion at the same layer.
\end{proof}

\section{The Downstream Chain Must Stay Fixed}

\begin{theorem}[Any admissible deeper continuation preserves the same downstream chain]
\label{thm:fixed}
Let \(\DeepSector\) be an admissible deeper continuation in the sense of Definition \ref{def:admissible}. Then the same recovered density/current block, the same composite primitive bridge-order block, the same phase-neutral minimizing branch, the same canonical profile \eqref{eq:canonicalprofile}, and the same dressed bridge map are recovered unchanged.
\end{theorem}

\begin{proof}
By Definition \ref{def:admissible}(2), the deeper continuation induces exactly the same raw-orbit lift \eqref{eq:rawliftA}--\eqref{eq:rawliftB} from the pair \((\ShapeTuple,\ScaleMod)\) to \(\RawTuple(\ScaleMod)\). By Definition \ref{def:admissible}(3), the minimizing vacuum branch remains phase-neutral, so the vacuum-angle sector introduces no new branch-selection mechanism or observer-visible angular forcing. By Definition \ref{def:admissible}(4), every downstream object recorded in \(\DownPkg\) depends on the deeper layer only through that same induced data.

Therefore no downstream formula changes. In particular, the ratio \(\RespRawEq/\RespRawRef\) remains equal to \(\RespShapeEq\), so the canonical profile reconstruction remains
\[
\CurvProfileCan(x)=1+\RespShapeEq\,x^2
=
1+\frac{\RespRawEq}{\RespRawRef}x^2.
\]
The same argument applies to the density/current block, the composite primitive bridge-order block, the phase-neutral minimizing branch, and the dressed bridge map: once the induced raw tuple and the mode of minimization are the same, the downstream outputs are the same.
\end{proof}

\begin{remark}
Theorem \ref{thm:fixed} is the exact reason the next note should not try to earn novelty by changing downstream formulas. At this stage of the line, the honest work is not to replace the successful downstream chain but to deepen its origin while keeping it intact.
\end{remark}

\section{Scale-Orbit and Vacuum-Angle Revisions Are Conditional}

\begin{definition}[Genuine quotient or locking data]
\label{def:quotient}
A deeper continuation carries \emph{genuine quotient or locking data} if it provides at least one of the following:
\begin{enumerate}
    \item a true identification of distinct positive moduli \(\ScaleMod^{(1)}\neq \ScaleMod^{(2)}\) as the same physical point of the deeper substrate theory;
    \item a dynamical or symmetry-forced mechanism selecting a distinguished modulus, thereby removing the present positive orbit family as physically distinct representatives;
    \item an anisotropic locking term or true quotient that identifies or fixes different constant vacuum angles \(\RespVacPhase=\phi_0\).
\end{enumerate}
\end{definition}

\begin{theorem}[Without genuine quotient or locking, the present verdicts remain]
\label{thm:status}
Let \(\DeepSector\) be an admissible deeper continuation. If \(\DeepSector\) does not supply genuine quotient or locking data in the sense of Definition \ref{def:quotient}, then:
\begin{enumerate}
    \item the positive scale orbit remains a canonical-normalization orbit rather than a proved deeper redundancy;
    \item the global \(O(2)\) vacuum angle remains residual.
\end{enumerate}
\end{theorem}

\begin{proof}
If distinct positive moduli are not identified, fixed, or quotiented by the deeper layer, then the family \(\ScaleMod>0\) remains a family of distinct representatives. The present canonical normalization is therefore still the disciplined way to choose one comparison slice, but it is not yet a proved redundancy. This gives item (1).

Likewise, if the deeper layer supplies neither an anisotropic locking term nor an angle quotient, then constant vacuum angles remain unforced and mutually distinct along the phase-neutral minimizing family. Hence the global \(O(2)\) angle stays residual. This gives item (2).
\end{proof}

\begin{corollary}[When revision is honest]
\label{cor:revision}
The scale-orbit verdict and the vacuum-angle verdict may be revised only on the basis of new quotient or locking data arising from the still deeper layer itself.
\end{corollary}

\begin{proof}
This is the contrapositive reading of Theorem \ref{thm:status}. Without genuine quotient or locking data, the current verdicts remain. Therefore an honest revision requires exactly such data.
\end{proof}

\begin{remark}
Corollary \ref{cor:revision} is deliberately conservative. It does not deny that a later note may eventually quotient the positive scale orbit or lock the angle. It says only that the raw-vacuum-orbit layer did not do so, and that the next note should not pretend otherwise unless the deeper substrate structure itself supplies the missing mechanism.
\end{remark}

\section{Claim Boundary}

This note does not claim that the still deeper orbit-shape/modulus origin has already been completed. It does not claim that the unit-reference shape block \(\ShapeTuple\) is already uniquely derived from first principles, and it does not claim that the positive scale orbit or the global \(O(2)\) vacuum angle have already been removed by a true quotient or locking mechanism.

Its claim is narrower. The project has already crossed the flagship exact-theory threshold, so the derivational bridge can now be handled with stronger discipline. At current scope the raw-vacuum-orbit note has already done the honest work at the raw-orbit layer. The remaining optional burden is deeper than that layer, and a successful continuation must be judged by the standards fixed here: derive or sharply constrain \(\ShapeTuple\) and \(\ScaleMod\), preserve the same downstream package \(\DownPkg\), and change the scale-orbit or vacuum-angle verdicts only if genuine quotient or locking data actually appear.

\section{Conclusion}

The repository now has a clean decision structure. If the priority is the flagship exact-theory statement of \TheoryName{}, the next work is packaging, integration, and presentation rather than another derivational manuscript. If the deeper derivational bridge is continued anyway, the next honest move is no longer a same-layer repair below the raw-vacuum-orbit note. It is a still deeper orbit-shape/modulus-origin, anchoring, or quotient manuscript.

This note fixes that burden precisely. The next admissible deeper manuscript must derive or sharply constrain the unit-reference orbit-shape block and the positive orbit modulus from still more primitive substrate structure, prove that the same downstream density/current block, primitive bridge-order block, phase-neutral minimizing branch, canonical profile, and dressed bridge map survive unchanged, and revisit the scale-orbit and vacuum-angle verdicts only if the deeper layer naturally supplies genuine quotient or locking structure. Until such a manuscript exists, the disciplined current verdict remains: the positive scale orbit is a canonical-normalization orbit rather than a proved redundancy, and the global \(O(2)\) vacuum angle remains residual.

\end{document}
